\section{Derive the exponential map and the Langer term}
\label{app:langer-map}
% BEGIN CURATED SUBJECT INDEX
\index{Langer correction}
\index{radial equation}
% END CURATED SUBJECT INDEX

Set $r=\rme^x$ and write
\begin{equation}
  u_l(r)=\rme^{x/2}\phi_l(x) .
  \label{eq:langer-wave-map}
\end{equation}
Using
\begin{equation}
  \frac{\rmd^2u_l}{\rmd r^2}
  =\rme^{-3x/2}
  \left(\frac{\rmd^2\phi_l}{\rmd x^2}
  -\frac{1}{4}\phi_l\right),
  \label{eq:langer-derivative}
\end{equation}
the radial equation becomes
\begin{equation}
  \left[
  -\frac{\hbar^2}{2m}\frac{\rmd^2}{\rmd x^2}
  +\rme^{2x}\bigl(V(\rme^x)-E\bigr)
  +\frac{\hbar^2}{2m}
  \left(l+\frac{1}{2}\right)^2
  \right]\phi_l(x)=0 .
  \label{eq:langer-transformed-equation}
\end{equation}
The $1/4$ shift is exact; it comes from removing the first derivative generated
by the coordinate change.  The origin condition has become a convergence
condition at $x\to-\infty$.  Equation~\eqref{eq:langer-transformed-equation}
therefore gives both an algebraic derivation of the Langer term and a geometric
derivation of the open-path/closed-cycle equivalence.
